\section{A quantitative first-layer spectrum bound}
\label{sec:critical-first-spectrum}

The proof starts with a quantitative lower bound on every Fourier
eigenvalue of the first tanh feature kernel.

\begin{lemma}[A quantitative first-layer spectrum bound]
\label{lem:first-tanh-spectrum}
Let $d\ge1$, $D=d+1$, $m=2^d$, and let $Q$ be the $m\times m$ matrix
\[
 Q(x,y)=\mathbb E_w\tanh(w\cdot\bar x)\tanh(w\cdot\bar y),
 \qquad w\sim\mathcal N(0,I_{d+1}/d).
\]
Then
\[
 \lambda_{\min}(Q)\ge e^{-13D}.
\]
\end{lemma}

\begin{proof}
Write $q=D/d\in[1,2]$, and let $h_n$ be the coefficient of the
normalized probabilists' Hermite polynomial $\operatorname{He}_n/\sqrt{n!}$
in $z\mapsto\tanh(\sqrt q z)$. The coefficients vanish for even $n$.
For every odd $n\ge1$, the Fourier sine identity
\[
 \tanh(\sqrt q z)=\frac1{\sqrt q}\int_0^\infty
       \frac{\sin(tz)}{\sinh(\pi t/(2\sqrt q))}\,dt
\]
and $\mathbb E[\operatorname{He}_n(Z)e^{itZ}]=(it)^ne^{-t^2/2}$ give
\[
 |h_n|=\frac1{\sqrt{q n!}}\int_0^\infty
       \frac{t^ne^{-t^2/2}}{\sinh(\pi t/(2\sqrt q))}\,dt.
\]
The sine identity follows, for example, by expanding
$1/\sinh(\pi t/2)=2\sum_{j\ge0}e^{-\pi(j+1/2)t}$ and using the
logarithmic derivative of the Euler product for $\cosh$.
Fubini is valid in the displayed Hermite integral: near zero use
$|\sin(tZ)|\le t|Z|$, and at infinity the reciprocal hyperbolic sine
has exponential decay.

Restrict the positive integral to $[\sqrt n,\sqrt n+1]$.
Using $n!\le n^n$, $1/\sinh u\ge2e^{-u}$, and $1\le q\le2$ gives
\[
 |h_n|\ge\sqrt2\exp\left[-\frac n2
               -\left(1+\frac\pi2\right)\sqrt n
               -\frac{1+\pi}{2}\right]\ge e^{-6n}.
\]
The last inequality follows from $\sqrt n\le n$, $n\ge1$ and $\pi<4$.

Mehler's identity now yields the nonnegative power series
\[
 Q(x,y)=\sum_{n\ \mathrm{odd}}h_n^2
       \left(\frac{1+\sum_{i=1}^d x_i y_i}{D}\right)^n.
\]
The Boolean characters diagonalize this translation-invariant kernel.
For a character of degree $r$, choose the odd integer $n=r$ if $r$
is odd and $n=r+1$ if $r$ is even. Then $1\le n\le D$.
In the multinomial expansion of the $n$th power, choosing every
selected coordinate exactly once, and choosing the constant coordinate
once when $r$ is even, contributes $n!/D^n$ to this character's
normalized eigenvalue. Every other contribution is nonnegative.
Therefore its eigenvalue for $Q/m$ is at least
\[
 h_n^2\frac{n!}{D^n}\ge e^{-12n-D}\ge e^{-13D}.
\]
Here $n!\ge(n/e)^n$, and
$n[1+\log(D/n)]\le D$, prove $n!/D^n\ge e^{-D}$.
Multiplying by $m\ge1$ proves the stated lower bound for $Q$.
\end{proof}

